5. 现在令 (8) 中的次数 $p,q$ 分别等于 $ld^{n-1}$ 和 $md^{n-1}$，即等于 (5) 中 $f(\xi)$ 和 $g(\xi)$ 的次数。把指标 $\mu,\nu$ 分拆为 $\mu^{(1)},\nu^{(1)},\mu^{(2)},\nu^{(2)}$；其中 $\mu^{(1)}$ 遍历 $f(\xi)$ 中缺失的所有指数，$\nu^{(1)}$ 遍历 $\psi(\xi,\eta)$ 中 $\xi$ 的所有指数；相应地，$\nu^{(2)}$ 遍历 $g(\xi)$ 中缺失的所有指数，$\mu^{(2)}$ 遍历 $\varphi(\xi,\eta)$ 中 $\xi$ 的所有指数。再令
\[
 e(\xi)=\sum Z_{k_1\ldots k_n}\xi^k\qquad(k_1+\cdots+k_n\le l+m)
\]
表示与多项式 (3) 相应的一元多项式；
\[
 S(Z,u)
\]
表示关于 $Z$ 和 $u$ 的整系数多项式；它由 $T(t,u)$ --- 即 (10) 中唯一确定的右端，其中 (9) 的求和遍及上述指标 $\mu^{(1)},\nu^{(1)},\mu^{(2)},\nu^{(2)}$ --- 通过把 $t$ 替换为系数 $Z$ 以及 $e(\xi)$ 的相应零点而得到。

现在把 $r,s$ 替换为 $f(\xi)$ 和 $g(\xi)$ 的系数 $A,B$ 以及相应零点，则在指定的求和下，经此代入后 $\zeta(u)$ \emph{消失}。此外，$t=w(r,s)$ 变为 $h(\xi)$ 的系数 $C(A,B)$ 及其相应零点；因此 $T(t,u)$ 变为 $S(C(A,B),u)$。由恒等式 (11) 便得到
\srcnumdisplay{(14)}{%
S(C(A,B),u)=0.\qquad[\text{关于 }A,B,u\text{ 恒等}].
}
由于 (14) 对每一个特化值系都保持成立，所以在一个\emph{扩张域中分解为两个因子}的这种特化 $\bar H(x)$ 或一般的 $\bar H(x,y)$，其系数 $\Gamma=C(A,B)$ 满足条件
\srcnumdisplay{(15)}{%
S(\Gamma,u)=0\qquad[\text{关于 }u\text{ 恒等}].
}

反过来，如果多项式 $\bar H(x)$，或形式 $\bar H(x,y)$ 的不全为零的系数 $\Gamma$ 满足条件 (15)，那么根据上面的论述，这等价于 $T(\tau,u)=0$，其中 $\tau$ 表示 $\bar h(\xi,\eta)$ 的全部系数。因而根据 (13)，在 $\Gamma$ 的一个扩张域中存在分解：
\[
 \bar h(\xi,\eta)=\bar f(\xi,\eta)\cdot\bar g(\xi,\eta),
\]
其中，由于指定的求和，$\bar f$ 和 $\bar g$ 中不出现任何不同于诱导指数的指数；于是存在关系
\[
 \bar H(x,y)=\bar F(x,y)\cdot\bar G(x,y).
\]
由此，只要在 $y=1$ 时没有因子变为零次，也就是说，特别是在 $\Gamma$ 不会导致 $\bar H(x)$ 次数降低时，还可得到关系
\[
 \bar H(x)=\bar F(x)\cdot\bar G(x).
\]

\emph{因此，条件 (15) 是 $\bar H(x,y)$ 在一个扩张域中分解为两个因子的充要条件；在排除次数降低时，它也是 $\bar H(x)$ 作这种分解的充要条件。}

由此还可推出，$S(Z,u)$ 不可能关于 $Z$ 和 $u$ 恒等消失，因为第 1 节已经证明，当 $n\ge2$ 时，以不定元为系数的多项式，因而也包括相应的多一个变量的齐次形式，都是绝对不可约的。因此，若用 $U$ 表示 $u$ 的幂积，则有：
\[
S(Z,u)=\sum \Phi_\lambda(Z)U_\lambda,
\]
其中 $\Phi_\lambda(Z)$ 根据 (14) 成为属于 $\mathfrak P$ 的多项式。\emph{由此证明了，$\mathfrak P$ 除了由参数表示 $\Gamma=C(A,B)$ 给出的零点以外，没有其他零点；这里 $A$ 和 $B$ 属于 $\Gamma$ 的一个代数扩张域。}这样便证明了第 2 节末尾所陈述的逆命题。

6. \emph{现在可以直接表述绝对不可约性的条件。}为此，只需把 (3) 中 $E(x)$ 的次数以所有可能的方式分解为两个正数之和 $l+m$，并为每个分解构造形式 $S(Z,u)$。这些形式的乘积
\srcnumdisplay{(16)}{%
R(Z,u)=\prod S(Z,u)
}
便构成 $E(x)$ 的\emph{可约性形式}；因为根据上面的论述，$\bar E(x)$ 或 $\bar E(x,y)$ 的每一次\emph{分解}，都对应于 (16) 的一个因子消失，反过来也成立。这样便证明了\emph{引言中提出的主定理}，同时也说明了如何用有限多个步骤实际构造可约性形式。

7. 由\emph{可约性形式的存在性} (16)，立即可以推出引言中提到的\emph{奥斯特罗夫斯基定理}。设
\[
 \bar E(x)=\sum \Gamma_{k_1\ldots k_n}x_1^{k_1}\cdots x_n^{k_n}
 \qquad(k_1+\cdots+k_n\le v)
\]
是一个以代数整数 $\Gamma$ 为系数的绝对不可约多项式；$\frK$ 表示包含所有 $\Gamma$ 的一个有限代数数域，$\mathfrak p$ 表示 $\frK$ 中的一个素理想。

那么，为 $\bar E(x)$ 的准确次数构造的可约性形式 $R(Z,u)$，在代入 $Z=\Gamma$ 后仍然\emph{不为零}；即
\[
 R(\Gamma,u)=\sum P_\lambda U_\lambda\ne0\qquad[\text{关于 }u\text{ 恒等}].
\]
因而，理想 $\mathfrak r=(P_1,P_2,\ldots)$ 只能被 $\frK$ 中\emph{有限多个}素理想整除。因此，如果 $\mathfrak p$ 不同于这有限多个素理想，那么 $R(\Gamma,u)$ 模 $\mathfrak p$ 不会恒等消失；又因为模 $\mathfrak p$ 的剩余类仍然构成一个域，所以可以推出，$\bar E(x)$ \emph{模 $\mathfrak p$ 不可约}，更准确地说，是绝对不可约的。$\bar E(x,y)$ 也是如此，因此模 $\mathfrak p$ 时同样不会发生次数降低。

如果 $\bar E(x)$ 的系数是代数分数，还必须要求 $\mathfrak p$ 不同于整除公分母 $\Delta$ 的有限多个素理想；模其余所有素理想时，可以用 $\Delta\cdot\bar E(x)$ 代替 $\bar E(x)$，从而满足上述假设。定理得证。

\begin{flushleft}
埃朗根，1921 年 8 月。
\end{flushleft}
\begin{center}
（1921 年 8 月 28 日收到。）
\end{center}

\clearpage
